% generated by roam.py handbook export -- не редактировать, править заметки
% 20261014225129-оценка_логарифма_сверху.org (body d8b4a5b821bf)

\begin{lemma}[{Оценка логарифма сверху}]\label[lemma]{note:41B4A714-C1F8-4267-9948-E13E99B30587}
  При $x > 0$
  \[
  \ln x \le x - 1
  \]
\end{lemma}

\begin{proof}
  Рассмотрим $- \ln x$ \\
  Это \hyperref[note:1CA14F2C-637F-41DE-9E97-72075D19E97F]{выпуклая функция} по \hyperref[note:86FCE16B-C82E-4416-851C-0E0AD0FDA382]{критерию выпуклости в терминах вторых производных}, так как
  \[
  \frac{\partial^2}{\partial x^2} (-\ln x) = \frac{1}{x^2} > 0
  \]

  Тогда по \hyperref[note:2F475EE4-2630-45FD-8DD2-8EA9D48059E4]{Дифференциальный критерий выпуклости функции одной переменной}
  \[
  f(x) := -\ln x \Rightarrow f(x) \ge f(1) + f^{\prime}(1) \cdot (x-1)
  \]
  Заметим, что $f(1) = 0$, $f^{\prime}(x) = - \frac{1}{x} \Rightarrow f^{\prime}(1) = -1$

  Тогда
  \[-
  \ln x \ge - (x - 1) \Leftrightarrow \ln x \le x - 1
  \]
\end{proof}
